\newpage
\section{Basic Typing and Well-Formedness Rules}\label{sec:tech:basic-typing}

\begin{figure}[ht!]
  {\small
  \begin{align*}
    \text{\textbf{Type erasure}} \quad & \eraserf{\ort{b}{\phi}} \doteq b \quad \eraserf{\urt{b}{\phi}} \doteq b \quad \eraserf{\tau_1 \ctinter \tau_2} \doteq \eraserf{\tau_1}
    \\&\eraserf{\tau_1 \ctoplus \tau_2} \doteq \eraserf{\tau_1} \quad  \eraserf{x{:}\tau_x \ctsarr \tau} \doteq \eraserf{\tau_x} \sarr \eraserf{\tau} \quad
    \eraserf{x{:}\tau_x \ctwand \tau} \doteq \eraserf{\tau_x} \sarr \eraserf{\tau} \quad
    \eraserf{\ctpers \tau} \doteq \eraserf{\tau}
    \\& \eraserf{\emptyset} \doteq \emptyset \quad \eraserf{x{:}\tau} \doteq x{:}\eraserf{\tau} \quad \eraserf{\Gamma_1 \ctcomma \Gamma_2} \doteq \eraserf{\Gamma_1}, \eraserf{\Gamma_2} \quad \eraserf{\Gamma_1 \ctstar \Gamma_2} \doteq \eraserf{\Gamma_1}, \eraserf{\Gamma_2}  \quad \eraserf{\Gamma_1 \ctoplus \Gamma_2} \doteq \eraserf{\Gamma_1}
 \end{align*}}
\caption{Type and Typing Context Erasure}
\label{fig:type-erasure}
\end{figure}

\begin{figure}[ht!]
    {\footnotesize
    {\small
   \begin{flalign*}
    &\text{\textbf{Basic Typing }} & \fbox{$\Sigma \basicvdash e : s$}
   \end{flalign*}}
   \begin{prooftree}
   \hypo{}
   \infer1[\textsc{BtConst}]{
   \Sigma \basicvdash c : \Code{BasicTy}(c)
   }
   \end{prooftree}
   \quad
   \begin{prooftree}
   \hypo{ \Sigma(x) = s }
   \infer1[\textsc{BtVar}]{
   \Sigma \basicvdash x : s
   }
   \end{prooftree}
   \quad
   \begin{prooftree}
   \hypo{\Sigma, x{:}s_x \basicvdash e : s}
   \infer1[\textsc{BtFun}]{
   \Sigma \basicvdash \lambda x{:}s_x.e : s_x{\sarr}s
   }
   \end{prooftree}
   \\[.8em]
   \begin{prooftree}
   \hypo{\Sigma, f{:}b\sarr s \basicvdash \lambda x{:}b.e : b\sarr s}
   \infer1[\textsc{BtFix}]{
   \Sigma \basicvdash \zfix{f}{b\sarr s}{x}{b}{e}
   : b{\sarr}s
   }
   \end{prooftree}
   \quad
   \begin{prooftree}
   \hypo{\Code{BasicTy}(\primop) = \overline{s_i}{\sarr}s \quad \forall i. \Sigma \basicvdash v_i : s_i}
   \infer1[\textsc{BtPrimOp}]{
   \Sigma \basicvdash \primop\, \overline{v_i} : s
   }
   \end{prooftree}
   \\[.8em]
   \begin{prooftree}
   \hypo{\Sigma \basicvdash v_1 : s_2 {\sarr} s \quad \Sigma \basicvdash v_2 : s_2}
   \infer1[\textsc{BtApp}]{
   \Sigma \basicvdash v_1\,v_2 : s
   }
   \end{prooftree}
   \\[.8em]
   \begin{prooftree}
   \hypo{\emptyset \basicvdash e_1 : s_x \quad \Sigma, x{:}s_x \basicvdash e_2 : s}
   \infer1[\textsc{BtLetE}]{
   \Sigma \basicvdash \zlet{x:s_x}{e_1}{e_2} : s
   }
   \end{prooftree}
   \quad
   \begin{prooftree}
   \hypo{\Sigma \basicvdash v : s_v \quad
   \forall i, \Code{BasicTy}(d_i) = \overline{s_j} {\sarr} s_v \quad
   \Sigma, \overline{y_j{:}s_j} \basicvdash e_i : s}
   \infer1[\textsc{BtMatch}]{
   \Sigma \basicvdash \match{v} \overline{d_i\,\overline{y_j} \to e_i} : s
   }
   \end{prooftree}
   }
   \caption{Basic Typing Rules}
       \label{fig:basic-type-rules}
   \end{figure}

   \begin{figure}[ht!]
    {\footnotesize
    {\small
   \begin{flalign*}
    &\text{\textbf{Basic Qualifier Typing }} & \fbox{$\Sigma \basicvdash \phi: s$}
   \end{flalign*}}
   \begin{prooftree}
   \hypo{\Code{BasicTy}(c) = s }
   \infer1[\textsc{BtConst}]{
   \Sigma \basicvdash c : s
   }
   \end{prooftree}
   \quad
   \begin{prooftree}
   \hypo{ \Sigma(x) = s }
   \infer1[\textsc{BtVar}]{
   \Sigma \basicvdash x : s
   }
   \end{prooftree}
   \quad
   \begin{prooftree}
   \hypo{ }
   \infer1[\textsc{BtTop}]{
   \Sigma \basicvdash \top : \Code{bool}
   }
   \end{prooftree}
   \quad
   \begin{prooftree}
   \hypo{ }
   \infer1[\textsc{BtBot}]{
   \Sigma \basicvdash \bot : \Code{bool}
   }
   \end{prooftree}
   \\[.8em]
   \begin{prooftree}
   \hypo{\Code{BasicTy}(\primop) = \overline{s_i}{\sarr}s \quad \forall i. \Sigma \basicvdash \phi_i : s_i }
   \infer1[\textsc{BtOp}]{
   \Sigma \basicvdash \primop\,\overline{\phi_i} : s
   }
   \end{prooftree}
   \quad
   \begin{prooftree}
   \hypo{
   \Sigma \basicvdash \phi : \Code{bool} }
   \infer1[\textsc{BtNeg}]{
   \Sigma \basicvdash \neg \phi : \Code{bool}
   }
   \end{prooftree}
   \\[.8em]
   \begin{prooftree}
   \hypo{
   \Sigma \basicvdash \phi_1 : \Code{bool} \quad
   \Sigma \basicvdash \phi_2 : \Code{bool} }
   \infer1[\textsc{BtAnd}]{
   \Sigma \basicvdash \phi_1 \land \phi_2 : \Code{bool}
   }
   \end{prooftree}
   \quad
   \begin{prooftree}
   \hypo{
   \Sigma \basicvdash \phi_1 : \Code{bool} \quad
   \Sigma \basicvdash \phi_2 : \Code{bool} }
   \infer1[\textsc{BtOr}]{
   \Sigma \basicvdash \phi_1 \lor \phi_2 : \Code{bool}
   }
   \end{prooftree}
   \quad
   \begin{prooftree}
   \hypo{
   \Sigma, x{:}b \basicvdash \phi : \Code{bool} }
   \infer1[\textsc{BtForall}]{
   \Sigma \basicvdash \forall x{:}b. \phi : \Code{bool}
   }
   \end{prooftree}
   }
   \caption{Basic Qualifier Typing Rules}
       \label{fig:basic-qualifier-type-rules}
   \end{figure}


\begin{figure}[ht!]
    {\footnotesize
    {\small
   \begin{flalign*}
    &\text{\textbf{Well-Formedness }} & \fbox{$\Sigma \wfvdash \tau$ \quad $\Sigma \wfvdash \Gamma$}
   \end{flalign*}}
   \begin{prooftree}
   \hypo{\Sigma, \vnu{:}b \basicvdash \phi : \Bool}
   \infer1[\textsc{WfBaseOver}]{
    \Sigma \wfvdash \ort{b}{\phi}
   }
   \end{prooftree}
   \quad
   \begin{prooftree}
    \hypo{\Sigma, \vnu{:}b \basicvdash \phi : \Bool}
    \infer1[\textsc{WfBaseUnder}]{
     \Sigma \wfvdash \urt{b}{\phi}
    }
    \end{prooftree}
    \\[.8em]
    \begin{prooftree}
        \hypo{\Sigma \wfvdash \tau_x \quad \Sigma, x{:}\eraserf{\tau_x} \wfvdash \tau}
        \infer1[\textsc{WfArr}]{
         \Sigma \wfvdash x{:}\tau_x \ctsarr \tau
        }
    \end{prooftree}
        \quad
    \begin{prooftree}
        \hypo{\Sigma \wfvdash \tau_x \quad \Sigma \wfvdash \tau}
        \infer1[\textsc{WfWand}]{
            \Sigma \wfvdash x{:}\tau_x \ctwand \tau
        }
    \end{prooftree}
    \quad
    \begin{prooftree}
        \hypo{\Sigma \wfvdash \tau}
        \infer1[\textsc{WfPers}]{
            \Sigma \wfvdash \ctpers \tau
        }
    \end{prooftree}
    \\[.8em]
    \begin{prooftree}
        \hypo{\Sigma \wfvdash \tau_1 \quad \Sigma \wfvdash \tau_2 \quad \eraserf{\tau_1} = \eraserf{\tau_2}}
        \infer1[\textsc{WfInter}]{
         \Sigma \wfvdash \tau_1 \interty \tau_2
        }
    \end{prooftree}
    \quad
    \begin{prooftree}
        \hypo{\Sigma \wfvdash \tau_1 \quad \Sigma \wfvdash \tau_2 \quad \eraserf{\tau_1} = \eraserf{\tau_2}}
        \infer1[\textsc{WfSum}]{
         \Sigma \wfvdash \tau_1 \ctoplus \tau_2
        }
    \end{prooftree}
    \\[.8em]
    \begin{prooftree}
        \hypo{}
        \infer1[\textsc{WfEmp}]{
            \Sigma \wfvdash \emptyset
        }
    \end{prooftree}
    \quad
        \begin{prooftree}
            \hypo{\Sigma \wfvdash \tau_x \quad x \notin \Dom{\Sigma}}
            \infer1[\textsc{WfSingle}]{
             \Sigma \wfvdash x{:}\tau_x
            }
        \end{prooftree}
    \\[.8em]
    \begin{prooftree}
        \hypo{\Sigma \wfvdash \Gamma_1 \quad \Sigma, \eraserf{\Gamma_1} \wfvdash \Gamma_2}
            \infer1[\textsc{WfComma}]{
             \Sigma \wfvdash \Gamma_1 \ctcomma \Gamma_2
        }
    \end{prooftree}
    \quad
    \begin{prooftree}
        \hypo{\Sigma \wfvdash \Gamma_1 \quad \Sigma \wfvdash \Gamma_2}
            \infer1[\textsc{WfStar}]{
             \Sigma \wfvdash \Gamma_1 \ctstar \Gamma_2
        }
    \end{prooftree}
    \\[.8em]
    \begin{prooftree}
        \hypo{\Sigma \wfvdash \Gamma_1 \quad \Sigma \wfvdash \Gamma_2 \quad \eraserf{\Gamma_1} = \eraserf{\Gamma_2}}
            \infer1[\textsc{WfCtxSum}]{
             \Sigma \wfvdash \Gamma_1 \ctoplus \Gamma_2
        }
    \end{prooftree}
   }
   \caption{Well-Formedness Rules for Context Types and Type Contexts}
       \label{fig:basic-refinement-type-rules}
   \end{figure}

   The type erasure function is defined in \autoref{fig:type-erasure}.
   The basic typing rules of our core language, qualifiers and
   refinement types are shown in \autoref{fig:basic-type-rules},
   \autoref{fig:basic-qualifier-type-rules}, and
   \autoref{fig:basic-refinement-type-rules}, resp. We use an
   auxiliary function $\Code{BasicTy}$ to provide a basic type for the
   primitives of our language, e.g., constants, built-in operators,
   and data constructors.
